\begin{workedexample}[Worked Example: Return loss to VSWR on a VNA]
A VNA reads $S_{11} = -14\ \text{dB}$ at band center. That magnitude is
$|\Gamma| = 10^{-14/20} = 0.20$, so
\[ \text{VSWR} = \frac{1 + 0.20}{1 - 0.20} = 1.5. \]
A short-open-load calibration first moves the reference plane to the probe tips,
so the reading describes the device rather than the cables.
\end{workedexample}

\begin{keytakeaways}
\begin{itemize}[leftmargin=1.4em,itemsep=2pt]
\item SOL (short-open-load) calibration sets the reference plane and removes systematic cable error.
\item $S_{11}$ in dB is return loss; convert to $|\Gamma|$ and VSWR as needed.
\item Port extension shifts the reference plane; time-domain gating isolates a reflection in space.
\end{itemize}
\end{keytakeaways}

\begin{exercises}
\item $S_{11} = -20\ \text{dB}$. What is the VSWR?
\item Which three standards define a basic one-port calibration?
\item Why calibrate at the probe tips rather than at the instrument port?
\end{exercises}

\furtherreading{Pozar~\cite{pozar} (ch.~4); Steer~\cite{steer}.}
